% Layout copied from the FINM 33165 DRL notes (itself modeled on the CAAM 31511 Monte Carlo notes):
% Times text with Computer Modern math, sans-serif headings, bold-italic
% proofs, ruled algorithm boxes, and rounded blue remark boxes.

\usepackage[T1]{fontenc}
\usepackage{times}
\usepackage{amsmath,amssymb,mathtools,bm}
\usepackage{amsthm,thmtools}
\usepackage[letterpaper,top=1.1in,bottom=1.1in,inner=1.15in,outer=1.15in,headheight=14pt]{geometry}
\usepackage{microtype}
\usepackage[shortlabels]{enumitem}
\usepackage{booktabs,array,tabularx}
\usepackage{xcolor}
\usepackage[most]{tcolorbox}
\usepackage{float}
\usepackage[font=small,labelfont=bf,textfont=it,labelsep=period]{caption}
\usepackage[chapter]{algorithm}
\usepackage{algpseudocode}
\usepackage{titlesec}
\usepackage{fancyhdr}
\usepackage{tikz}
\usetikzlibrary{arrows.meta,positioning,calc,fit}
\usepackage[colorlinks=true,linkcolor=blue!65!black,citecolor=blue!65!black,urlcolor=blue!65!black]{hyperref}
\usepackage[capitalise,nameinlink]{cleveref}

\definecolor{boxblue}{RGB}{30,60,170}

% ---------------------------------------------------------------- headings
\newcommand{\headsquare}{\raisebox{0.12ex}{\scalebox{0.55}{$\blacksquare$}}}
\titleformat{\chapter}[display]{\sffamily\bfseries}{\Large\chaptertitlename\ \thechapter}{8pt}{\Huge}
\titlespacing*{\chapter}{0pt}{40pt}{36pt}
\titleformat{\section}{\sffamily\bfseries\Large}{\thesection\hspace{0.35em}\headsquare}{0.35em}{}
\titleformat{\subsection}{\sffamily\bfseries\large}{\thesubsection\hspace{0.3em}\headsquare}{0.3em}{}
\titleformat{\subsubsection}{\sffamily\bfseries}{}{0pt}{}
\setcounter{secnumdepth}{2}
\setcounter{tocdepth}{1}
\setlength{\emergencystretch}{2em}

\pagestyle{fancy}
\fancyhf{}
\fancyhead[LE,RO]{\thepage}
\fancyhead[RE]{\leftmark}
\fancyhead[LO]{\rightmark}
\renewcommand{\headrulewidth}{0.4pt}
\renewcommand{\chaptermark}[1]{\markboth{Chapter \thechapter.\ #1}{}}
\renewcommand{\sectionmark}[1]{\markright{\thesection.\ #1}}
\fancypagestyle{plain}{\fancyhf{}\fancyfoot[C]{\thepage}\renewcommand{\headrulewidth}{0pt}}

% ---------------------------------------------------------------- theorems
\declaretheoremstyle[
  spaceabove=8pt, spacebelow=8pt,
  headfont=\bfseries, notefont=\bfseries, notebraces={(}{)},
  bodyfont=\itshape, headpunct={.}, postheadspace=0.5em
]{notesthm}
\declaretheoremstyle[
  spaceabove=8pt, spacebelow=8pt,
  headfont=\bfseries, notefont=\bfseries, notebraces={(}{)},
  bodyfont=\normalfont, headpunct={.}, postheadspace=0.5em
]{notesremark}
\declaretheorem[style=notesthm,numberwithin=chapter,name=Theorem]{theorem}
\declaretheorem[style=notesthm,sibling=theorem,name=Lemma]{lemma}
\declaretheorem[style=notesthm,sibling=theorem,name=Proposition]{proposition}
\declaretheorem[style=notesthm,sibling=theorem,name=Corollary]{corollary}
\declaretheorem[style=notesthm,sibling=theorem,name=Definition,qed=$\square$]{definition}
\declaretheorem[style=notesthm,sibling=theorem,name=Assumption,qed=$\square$]{assumption}
\declaretheorem[style=notesremark,sibling=theorem,name=Remark,qed=$\square$]{remark}

\renewcommand{\proofname}{\bfseries Proof}
\newenvironment{proofsketch}{\begin{proof}[\bfseries Proof sketch]}{\end{proof}}

\crefname{algorithm}{Algorithm}{Algorithms}
\crefname{assumption}{Assumption}{Assumptions}

% ---------------------------------------------------------------- algorithms
\algrenewcommand\algorithmicrequire{\textbf{Input:}}
\algrenewcommand\algorithmicensure{\textbf{Output:}}
\newcommand{\Input}[1]{\Statex \hspace*{-\algorithmicindent}\hspace*{1.2em}\textbf{Input:} #1}
\newcommand{\Init}[1]{\Statex \hspace*{-\algorithmicindent}\hspace*{1.2em}\textbf{Initialize:} #1}
\newcommand{\Output}[1]{\Statex \hspace*{-\algorithmicindent}\hspace*{1.2em}\textbf{Output:} #1}
\algrenewcommand\alglinenumber[1]{\footnotesize #1:}

% ---------------------------------------------------------------- boxes
% \begin{notebox}{Pros and Cons} ... \end{notebox}
\newtcolorbox{notebox}[1]{enhanced, colback=white, colframe=boxblue, boxrule=0.8pt,
  arc=3mm, left=6pt, right=6pt, top=4pt, bottom=4pt, before skip=10pt, after skip=10pt,
  before upper={\textcolor{boxblue}{\textbf{#1}:}\ }}

% ---------------------------------------------------------------- notation
\newcommand{\R}{\mathbb{R}}
\newcommand{\N}{\mathbb{N}}
\newcommand{\E}{\mathbb{E}}
\newcommand{\Prob}{\mathbb{P}}
\newcommand{\V}{\mathbb{V}}
\newcommand{\cS}{\mathcal{S}}
\newcommand{\cA}{\mathcal{A}}
\newcommand{\cF}{\mathcal{F}}
\newcommand{\cD}{\mathcal{D}}
\newcommand{\cH}{\mathcal{H}}
\newcommand{\cL}{\mathcal{L}}
\newcommand{\cN}{\mathcal{N}}
\newcommand{\cB}{\mathcal{B}}
\newcommand{\ind}{\mathbf{1}}
\newcommand{\KL}{D_{\mathrm{KL}}}
\newcommand{\given}{\,|\,}
\newcommand{\dd}{\,\mathrm{d}}
\DeclareMathOperator*{\argmax}{arg\,max}
\DeclareMathOperator*{\argmin}{arg\,min}
\DeclareMathOperator{\clip}{clip}
\DeclareMathOperator{\tr}{tr}
\newcommand{\iidsim}{\stackrel{\mathrm{i.i.d.}}{\sim}}
\newcommand{\norm}[1]{\lVert #1\rVert}
\newcommand{\supnorm}[1]{\lVert #1\rVert_\infty}

% ---------------------------------------------------------------- additions for the STA 314/414 notes
\DeclareMathOperator{\Cov}{Cov}
\DeclareMathOperator{\Var}{Var}
\DeclareMathOperator{\diag}{diag}
\DeclareMathOperator{\sign}{sign}
\DeclareMathOperator{\softmax}{softmax}
\DeclareMathOperator{\ELBO}{ELBO}
\DeclareMathOperator{\MSE}{MSE}
\DeclareMathOperator{\Err}{Err}
\newcommand{\cX}{\mathcal{X}}
\newcommand{\cY}{\mathcal{Y}}
\newcommand{\cZ}{\mathcal{Z}}
\newcommand{\cG}{\mathcal{G}}
\newcommand{\cQ}{\mathcal{Q}}
\newcommand{\cC}{\mathcal{C}}
\newcommand{\cT}{\mathcal{T}}
\newcommand{\cK}{\mathcal{K}}
\newcommand{\GP}{\mathcal{GP}}
\newcommand{\T}{^{\top}}
\newcommand{\hatf}{\hat f}
\newcommand{\indep}{\perp\!\!\!\perp}
\crefname{part}{Part}{Parts}
